\documentclass[12pt]{article}

% Core math
\usepackage{amsmath, amssymb, amsthm}
\usepackage{mathtools}
\usepackage{stmaryrd}
\usepackage{booktabs}

% Diagrams
\usepackage{tikz-cd}
\usepackage{tikz}

% References
\usepackage{hyperref}
\usepackage{cleveref}

% Graphics
\usepackage{graphicx}

% Page layout
\usepackage[margin=1in]{geometry}

% GrokRxiv sidebar
\usepackage{xcolor}

% Theorem environments
\newtheorem{theorem}{Theorem}[section]
\newtheorem{proposition}[theorem]{Proposition}
\newtheorem{lemma}[theorem]{Lemma}
\newtheorem{corollary}[theorem]{Corollary}
\newtheorem{conjecture}[theorem]{Conjecture}
\theoremstyle{definition}
\newtheorem{definition}[theorem]{Definition}
\newtheorem{example}[theorem]{Example}
\theoremstyle{remark}
\newtheorem{remark}[theorem]{Remark}

% Provide \emergencystretch for line-breaking in long sentences
\setlength{\emergencystretch}{2em}

% GrokRxiv sidebar definition (uses modern LaTeX hook)
\definecolor{grokgray}{RGB}{110,110,110}
\AddToHook{shipout/foreground}{%
  \ifnum\value{page}=1
    \begin{tikzpicture}[remember picture, overlay]
      \node[
        rotate=90,
        anchor=south,
        font=\Large\sffamily\bfseries\color{grokgray},
        inner sep=0pt
      ] at ([xshift=38pt, yshift=0.52\paperheight]current page.south west)
      {GrokRxiv:2026.05.etcs-izf-folds\quad
       [\,math.LO\,]\quad
       04 May 2026};
    \end{tikzpicture}
  \fi
}

\title{ETCS, IZF, and FOLDS:\\
Comparative Structural Foundations and the Univalence Boundary}

\author{Matthew Long \\
\textit{The YonedaAI Collaboration} \\
\textit{YonedaAI Research Collective} \\
Chicago, IL \\
\texttt{research@yonedaai.com} $\cdot$ \url{https://yonedaai.com}}

\date{4 May 2026}

\begin{document}

\maketitle

\begin{abstract}
We undertake a three-way structural comparison of three foundational systems: Lawvere's
Elementary Theory of the Category of Sets (ETCS, 1964), Friedman's Intuitionistic
Zermelo--Fraenkel set theory (IZF), and Makkai's First-Order Logic with Dependent Sorts
(FOLDS, 1995); and we measure each against the Homotopy Type Theory (HoTT) frame in
which the Univalence Axiom acts as a structural watermark. We give precise statements of
all three axiomatisations, present the bi-interpretation theorems linking them
(McLarty 2004 for ETCS $\equiv$ Bounded Zermelo + Replacement; Friedman--Aczel for IZF;
Makkai's invariance theorem for FOLDS), and analyse which foundational principles in each
system become \emph{theorems} (rather than axioms) in HoTT once Univalence is assumed.
We identify what we call the \emph{univalence boundary}: the locus separating those
structural principles whose validity in a foundation requires univalence (e.g.\ the
Structure Identity Principle, full FOLDS-equivalence, indiscernibility of structurally
isomorphic objects) from those that hold in much weaker fragments (e.g.\ existence of
NNO, Replacement, the validity of equivalence-invariant first-order proofs). We
formalise the comparison as a square of interpretations, identify three open problems on
the univalence side, and provide companion code in Haskell (with QuickCheck-style
invariance tests) and Lean~4 (Mathlib's \texttt{CategoryTheory}) demonstrating ETCS
axioms, FOLDS signatures, and the invariance check. The paper continues the categorical
foundations programme of Paper~VI of our prior series and complements Papers~II--V on
HoTT.
\end{abstract}

\tableofcontents

\section{Introduction}
\label{sec:intro}

\subsection{Why three foundations}

The phrase ``foundations of mathematics'' is in 2026 plural by both fact and conviction.
A working mathematician encountering structures up to isomorphism (vector spaces, groups,
topological spaces, schemes, $\infty$-categories) needs an underlying logic that respects
the equivalence principle: \emph{isomorphic objects share their mathematical content}.
Classical Zermelo--Fraenkel set theory with Choice (ZFC), the de facto standard since the
1920s, supports this principle only by a sociological convention --- the convention to
prove only structural properties --- but the formalism itself permits non-structural
predicates such as ``$\emptyset \in 7$'' (a sentence whose truth value depends on
Kuratowski's particular encoding of pairs).

The three foundational alternatives we examine in this paper each respond to that
inadequacy in a different way:
\begin{itemize}
\item \textbf{ETCS (Lawvere 1964)} replaces the universe of sets-with-membership by a
\emph{category} of sets-with-functions, axiomatised structurally as a well-pointed topos
with a natural numbers object (NNO) and a choice operator;
\item \textbf{IZF (Friedman 1973, Myhill, Aczel)} keeps the membership predicate but
moves to intuitionistic logic, accepting full Separation, Replacement, Powerset, and
Infinity while rejecting Excluded Middle and Choice; structuralism is enforced indirectly
via realisability and topos-valued models;
\item \textbf{FOLDS (Makkai 1995)} re-engineers the underlying \emph{logic} so that only
isomorphism-invariant predicates are expressible; the equivalence principle is a metatheorem
of the syntax rather than a moral imperative.
\end{itemize}

Each is motivated by the search for a logical environment in which equivalence-invariance
is built in. Each succeeds for a different reason. And each has a precise relationship
to HoTT, the foundation introduced by Voevodsky (2010) in which equivalence-invariance is
proven as a theorem (the Univalence Axiom and its corollary, the Structure Identity
Principle).

\subsection{Structural vs material set theory}

Following Awodey~\cite{Awodey2014}, McLarty~\cite{McLarty2004}, and Shulman, we use the
following terminology.

\begin{definition}[Material set theory]
A set theory is \emph{material} if its primitive non-logical relation is membership
$\in$, with sets identified by extensionality (two sets are equal iff they have the same
elements), and where the universe of discourse is a single class of all sets ranging over
a $\in$-tree.
\end{definition}

\begin{definition}[Structural set theory]
A set theory is \emph{structural} if its primitive notions are sets and functions
(equivalently: the objects and morphisms of a category), with sets identified up to
unique isomorphism, and where membership is a derived notion ($x\in A$ means
``$x:1\to A$ is a global element'').
\end{definition}

ZFC and IZF are material; ETCS and Lawvere's ETCC (Elementary Theory of the Category of
Categories) are structural; FOLDS is neither --- it is a logic, applicable to either
flavour, but designed so that even when interpreted materially, only structural sentences
are syntactically expressible. HoTT is structural (in fact \emph{higher-structural}: types
are identified up to higher equivalence).

\subsection{The univalence boundary}

The methodological question that organises this paper is:
\begin{quote}
\emph{Of the principles, axioms, and metatheorems that distinguish ETCS, IZF, and FOLDS,
which become theorems in HoTT (with univalence) and which do not?}
\end{quote}

A principle that becomes a theorem under HoTT+UA but not under HoTT alone (i.e.\ the
$(-1)$-truncated fragment, or the propositions-as-subsingletons fragment) lies on the
\emph{far} side of the univalence boundary: it is \emph{conditional on} univalence.
A principle that holds without univalence lies on the \emph{near} side: it is
unconditional.

We will argue in \cref{sec:boundary} that the boundary cuts as follows:

\begin{center}\small
\begin{tabular}{lcccc}
\toprule
\textbf{Principle} & ETCS & IZF & FOLDS & UA in HoTT? \\
\midrule
NNO existence & axiom & thm & expr.\ & no \\
Function ext. & thm & thm & axiom & implied by UA\(^{*}\) \\
Prop.\ ext. & thm & thm & expr.\ & no \\
SIP & meta & --- & meta & \textbf{yes} \\
FOLDS-eq.\ $=$ id & --- & --- & meta & \textbf{yes} \\
EM & axiom & reject & expr.\ & indep. \\
Choice & AC & reject & expr.\ & indep.\ (PEM+AC) \\
Replacement & axiom & axiom & --- & thm (small) \\
Powerset & derived & axiom & --- & impred.\ $\Omega$ \\
\bottomrule
\end{tabular}
\end{center}

\noindent\(^{*}\)Function extensionality (FE) is a theorem of HoTT~+~UA (Voevodsky) but is
independent of MLTT alone; see \cref{sec:hott} and \cref{thm:fefar} for the precise
nuance.

\subsection{Outline}
\Cref{sec:etcs} states the eight-axiom presentation of ETCS, proves the elementary
consequences (NNO, products, exponentials), and recalls McLarty's bi-interpretation
theorem with Bounded Zermelo + Replacement.
\Cref{sec:izf} presents IZF, contrasts it with CZF and ZFC, and discusses the
realisability and topos-valued models that connect IZF to HoTT.
\Cref{sec:folds} develops Makkai's signature theory, defines FOLDS-equivalence, and
proves Makkai's Invariance Theorem (every FOLDS-formula is invariant under FOLDS-equivalence).
\Cref{sec:hott} reframes HoTT as the unifying framework: which structural principles
become theorems, which need univalence.
\Cref{sec:models} surveys models: ETCS in toposes, FOLDS in fibrations, HoTT in
$(\infty,1)$-toposes.
\Cref{sec:boundary} formally analyses the univalence boundary.
\Cref{sec:open} states three open problems.
\Cref{sec:conclusion} concludes.

\section{ETCS: The Elementary Theory of the Category of Sets}
\label{sec:etcs}

\subsection{The eight axioms}

ETCS is a first-order theory in the language of categories (objects, morphisms,
composition, identity). The eight axioms below follow Lawvere's original 1964 PNAS paper~\cite{Lawvere1964} as updated by Lawvere--Rosebrugh~\cite{LawvereRosebrugh2003} and
McLarty~\cite{McLarty2004}.

\begin{definition}[ETCS axioms]
\label{def:etcs}
A category $\mathcal{S}$ \emph{models ETCS} iff it satisfies:
\begin{enumerate}
\item[(A1)] $\mathcal{S}$ has finite limits (terminal object $1$, binary products $A\times B$,
equalisers).
\item[(A2)] $\mathcal{S}$ is Cartesian closed: for any objects $A,B$ there is an
exponential $B^A$ with evaluation $\mathrm{ev}: B^A \times A \to B$ enjoying the universal
property of $\lambda$-abstraction.
\item[(A3)] $\mathcal{S}$ has a subobject classifier: an object $\Omega$ with a morphism
$\top: 1 \to \Omega$ such that every monomorphism $m: A \hookrightarrow B$ arises as the
pullback of $\top$ along a unique characteristic morphism $\chi_m : B \to \Omega$.
\item[(A4)] $\mathcal{S}$ has a natural numbers object: an object $\mathbb{N}$ with maps
$0:1\to \mathbb{N}$ and $s:\mathbb{N}\to \mathbb{N}$ such that for every $(X, x_0:1\to X,
f:X\to X)$ there is a unique morphism $h:\mathbb{N}\to X$ with $h\circ 0 = x_0$ and
$h\circ s = f\circ h$.
\item[(A5)] $\mathcal{S}$ is well-pointed: if $f,g:A\to B$ are distinct morphisms, then
there exists $x:1\to A$ with $f\circ x \ne g\circ x$.
\item[(A6)] (Axiom of choice) Every epimorphism $e: A \twoheadrightarrow B$ has a section
$s: B \to A$ with $e \circ s = \mathrm{id}_B$.
\item[(A7)] (Two-valuedness; derivable from A5--A6 and A8) Any morphism $f:1\to \Omega$
equals either $\top$ or $\bot$, where $\bot$ is the unique map factoring through the
initial object. Listed for didactic completeness; see the remark below.
\item[(A8)] (Non-degeneracy) $0\not\cong 1$ (the initial object is not isomorphic to the
terminal).
\end{enumerate}
\end{definition}

\begin{remark}
A1+A2+A3 alone make $\mathcal{S}$ an \emph{elementary topos} in the sense of
Lawvere--Tierney. Adding A4 gives a \emph{topos with NNO}. Adding A5--A8 specialises to
``topos behaving like classical sets''. Equivalently, an ETCS model is a well-pointed
topos with NNO and choice. Two-valuedness (A7) follows from well-pointedness plus
non-degeneracy when AC is present.
\end{remark}

\subsection{Elementary consequences}

\begin{proposition}[Function extensionality in ETCS]
\label{prop:etcsfunext}
For $f,g: A\to B$, if $\forall x:1\to A,\, f\circ x = g\circ x$, then $f=g$.
\end{proposition}
\begin{proof}
This is the contrapositive of well-pointedness (A5).
\end{proof}

\begin{proposition}[Power objects]
\label{prop:etcspower}
$\mathcal{S}$ has power objects $P(A) := \Omega^A$, with the standard adjunction
$\mathrm{Sub}(B\times A)\cong \mathrm{Hom}(B, P(A))$.
\end{proposition}
\begin{proof}
Combine A2 and A3: $P(A) = \Omega^A$ classifies subobjects of any product
$B\times A$ via $\chi: B\times A \to \Omega$ uncurried as $\widetilde{\chi}:
B\to \Omega^A$.
\end{proof}

\begin{theorem}[McLarty's bi-interpretation theorem]
\label{thm:mclarty}
ETCS and Bounded Zermelo set theory plus Replacement (\textbf{BZ}+\textbf{R}) are
bi-interpretable: each can interpret the other, and the interpretations compose to the
identity up to definable isomorphism.
\end{theorem}
\begin{proof}[Sketch]
Given a model of ETCS, define a class of \emph{trees} as well-founded extensional
quotients of pointed graphs internalised by the NNO and powerset; this gives a model of
$\mathbf{BZ}+\mathbf{R}$. Conversely, given a model $V$ of $\mathbf{BZ}+\mathbf{R}$,
define $\mathcal{S}_V$ to be the category of sets and functions inside $V$; one verifies
A1--A8 directly. Composition produces models that are equivalent
in the appropriate 2-categorical sense. Full proof in McLarty~\cite{McLarty2004}; we
return to it in \cref{sec:models}.
\end{proof}

\subsection{Encoding-free arithmetic}

ETCS resolves Benacerraf's identification problem (Paper~II of our series) by making
\emph{ordered pair} a primitive operation (the universal property of products) rather
than the Kuratowski encoding $\langle a,b\rangle = \{\{a\},\{a,b\}\}$. Junk theorems
such as ``$3 \in 7$'' (vacuously true under von Neumann ordinals: $3 = \{0,1,2\}$ and
$7=\{0,\ldots,6\}$) are not formulable in the ETCS language: $3:1\to \mathbb{N}$ is a
global element, $7:1\to \mathbb{N}$ likewise, and there is no membership predicate
between them.

\begin{example}
$\mathbb{N} \times \mathbb{N}$ in ETCS is the categorical product of $\mathbb{N}$ with
itself, characterised by the universal property; whether one uses Kuratowski pairs,
Wiener pairs, or any other internal encoding to \emph{prove existence} inside
$\mathbf{BZ+R}$ is irrelevant to the structural object so produced.
\end{example}

\subsection{Worked example: arithmetic via the recursor}

We illustrate how arithmetic is developed structurally in ETCS using only the universal
property of the NNO. Given the NNO $(\mathbb{N},0,s)$, define addition
$+:\mathbb{N}\times\mathbb{N}\to\mathbb{N}$ as follows.

\begin{definition}[Addition in ETCS]
\label{def:etcsadd}
The addition map $+: \mathbb{N}\times\mathbb{N}\to\mathbb{N}$ is the unique morphism
satisfying:
\begin{itemize}
\item $+ \circ \langle 0_{\mathbb{N}}, \mathrm{id}_{\mathbb{N}}\rangle = \mathrm{id}_{\mathbb{N}}$ (left zero), and
\item $+ \circ (s\times \mathrm{id}_{\mathbb{N}}) = s \circ +$ (left successor).
\end{itemize}
\end{definition}

\begin{proposition}[Existence of addition]
\label{prop:etcsadd}
There exists a unique morphism $+$ as in \cref{def:etcsadd}.
\end{proposition}
\begin{proof}
Apply the NNO universal property to the pointed endomorphism
$(\mathbb{N}^{\mathbb{N}}, \mathrm{id}_{\mathbb{N}} : 1\to \mathbb{N}^{\mathbb{N}},
\widehat{s\circ\mathrm{ev}}: \mathbb{N}^{\mathbb{N}}\to \mathbb{N}^{\mathbb{N}})$ where
the basepoint is the constant function $\mathrm{id}_{\mathbb{N}}$ and the endomap is
post-composition with $s$. This yields a unique
$\widehat{+}: \mathbb{N} \to \mathbb{N}^{\mathbb{N}}$, which uncurries to
$+: \mathbb{N}\times\mathbb{N}\to \mathbb{N}$ satisfying both equations.
\end{proof}

\begin{example}[Multiplication and exponentiation]
Multiplication $\cdot$ is defined by the same recursion technique with
$+$ in place of $s$:
$\cdot\circ\langle 0,\mathrm{id}\rangle = 0$, $\cdot\circ (s\times \mathrm{id}) =
+\circ\langle\mathrm{id},\cdot\rangle$.
Exponentiation $\uparrow$ similarly: $\uparrow\circ\langle 0,\mathrm{id}\rangle = 1$,
$\uparrow\circ (s\times \mathrm{id}) = \cdot\circ\langle\mathrm{id},\uparrow\rangle$.
\end{example}

\subsection{ETCS without choice}

Removing A6 (choice) yields what is sometimes called \emph{ETCS-AC} or \emph{Topos+NNO+WP}
(well-pointed topos with NNO). This system is consistent with both EM and $\neg$EM in
non-trivial models. McLarty's bi-interpretation theorem extends to: ETCS-AC is
bi-interpretable with $\mathbf{BZ}+\mathbf{R}$-AC. The constructive variant of ETCS
(replace A5 well-pointedness by Kock's notion of subterminal-pointedness, drop A6 and A7)
is bi-interpretable with appropriate fragments of IZF. We will see in
\cref{sec:boundary} that AC is on the \emph{straddling} side of the univalence boundary.

\subsection{Comparison with Awodey--Forssell categorical structuralism}

Awodey~\cite{Awodey2014} introduces a refined notion of \emph{structural foundation}: a
formal system whose models satisfy isomorphism-invariance internally. ETCS satisfies
this in the 1-categorical sense; its only non-invariant data are the choices of products,
exponentials, etc., which are determined up to (unique) isomorphism. The
``isomorphism-invariance'' guarantee in ETCS, however, is not a theorem of ETCS
\emph{about itself}; it is a metatheorem (any two products are isomorphic; any property
expressible in the categorical language is invariant under isomorphism). To make the
guarantee \emph{internal}, one needs HoTT.

\section{IZF: Intuitionistic Zermelo--Fraenkel}
\label{sec:izf}

\subsection{Axioms}

IZF (Friedman 1973, Myhill, surveyed in~\cite{StanfordIZFCZF}) is the system over intuitionistic first-order logic with equality
and a binary predicate $\in$, having the following axioms.

\begin{definition}[IZF axioms]
\label{def:izf}
We adopt the formulation of IZF using the strong \emph{Collection schema}
(rather than the weaker Replacement schema); under intuitionistic logic, Collection is
strictly stronger than Replacement, so this choice is significant. The remark following
the definition discusses the difference.
\begin{enumerate}
\item[(I1)] (Extensionality) $\forall x\forall y\,[(\forall z\,(z\in x \Leftrightarrow
z\in y))\Rightarrow x=y]$.
\item[(I2)] (Empty set) $\exists x\, \forall y\,(y\notin x)$.
\item[(I3)] (Pairing) $\forall x\forall y\, \exists z\,\forall w\,
(w\in z\Leftrightarrow w=x\vee w=y)$.
\item[(I4)] (Union) $\forall x\, \exists y\,\forall z\,(z\in y \Leftrightarrow \exists w
\in x.\, z\in w)$.
\item[(I5)] (Infinity) $\exists \omega\,(\emptyset\in\omega \wedge \forall x\in\omega.\,
x\cup\{x\}\in\omega)$.
\item[(I6)] (Full Separation) For every formula $\varphi(z)$, $\forall x\,\exists y\,
\forall z\,(z\in y \Leftrightarrow z\in x \wedge \varphi(z))$.
\item[(I7)] (Powerset) $\forall x\, \exists y\, \forall z\,(z\in y \Leftrightarrow z
\subseteq x)$.
\item[(I8)] (Collection) For every formula $\varphi(x,y)$, if $\forall x\in a.\,\exists
y.\, \varphi(x,y)$, then $\exists b.\,\forall x\in a.\,\exists y\in b.\, \varphi(x,y)$.
\item[(I9)] ($\in$-Induction) For every $\varphi$,
$[\forall x\,((\forall y\in x.\,\varphi(y))\Rightarrow \varphi(x))]\Rightarrow \forall x.\,
\varphi(x)$.
\end{enumerate}
The underlying logic is intuitionistic; \emph{Excluded Middle (EM)} and \emph{Choice (AC)}
are not assumed.
\end{definition}

\begin{remark}
A standard simplification replaces (I8) Collection by the weaker Replacement schema.
Friedman's IZF uses Collection because, under intuitionistic logic, Collection is
strictly stronger than Replacement: Replacement requires uniqueness of the witness, while
Collection requires only existence. In classical logic, Collection is derivable from
Replacement using EM; in IZF this implication fails.
\end{remark}

\subsection{IZF and ZFC: strength and divergence}

\begin{theorem}[Friedman~\cite{Friedman1973}]
\label{thm:friedman}
$\mathrm{ZF} = \mathrm{IZF} + \mathrm{EM}$, where $\mathrm{EM}$ is the schema
$\varphi\vee\neg\varphi$ over the IZF language. Consequently, IZF and ZF are equiconsistent.
\end{theorem}
\begin{proof}[Sketch]
Forward direction: $\mathrm{IZF}\subseteq \mathrm{ZF}$ definitionally, since classical
logic includes intuitionistic logic and EM gives the rest. Reverse: every ZF axiom is in
IZF except those needing EM in their classical forms, but Separation in IZF is full so
the classical schema follows. The equiconsistency: any model of ZF is a model of
IZF+EM; any model of IZF gives a Heyting-valued model that, under EM, becomes a model of
ZF. Friedman's original proof uses double-negation translations.
\end{proof}

\begin{remark}[CZF vs IZF]
Aczel's CZF (Constructive Zermelo--Fraenkel) replaces Powerset by \emph{Subset Collection}
(or its equivalent \emph{Fullness} axiom) and Full Separation by \emph{Restricted
Separation} (where $\varphi$ is bounded). CZF is significantly weaker than IZF:
$\mathrm{CZF} \subsetneq \mathrm{IZF}$, and CZF is interpretable in Martin-L\"of type
theory via Aczel's sets-in-types translation. IZF is impredicative
(Powerset and Full Separation), which puts it in tension with the predicative spirit of
MLTT and HoTT.
\end{remark}

\subsection{IZF, realisability, and HoTT}

\begin{definition}[Realisability model]
A \emph{realisability model} of IZF is a category-theoretic structure
$(\mathbb{V}, \Vdash)$ where $\mathbb{V}$ is a class hierarchy and $\Vdash$ is a forcing
relation specified by terms of an underlying calculus (e.g.\ untyped $\lambda$-calculus,
PCAs, or partial combinatory algebras) such that the IZF axioms are forced.
\end{definition}

The connection to HoTT runs through the topos-theoretic interpretation. Hyland's
effective topos $\mathcal{E}\!f\!f$ provides a realisability model in which IZF holds
(via Full Separation, Collection, and Powerset using the lifted $\Omega$). Inside
$\mathcal{E}\!f\!f$, the internal language is intuitionistic higher-order logic, which
embeds into HoTT once we resolve coherence at higher levels. Shulman's results~\cite{Shulman2019} show every Grothendieck $(\infty,1)$-topos models HoTT; specialising
to localic $(\infty,1)$-toposes gives realisability models of HoTT compatible with IZF.

\begin{theorem}[Awodey--Warren propositions-as-types for IZF]
\label{thm:awodey-warren}
The $(-1)$-truncated fragment of HoTT (i.e.\ HoTT restricted to mere propositions)
satisfies the IZF axioms in the following sense: there is an interpretation
$\llbracket\cdot\rrbracket$ of IZF formulas into $(-1)$-truncated types such that
$\mathrm{IZF}\vdash \varphi$ implies $\mathrm{HoTT}\vdash \llbracket\varphi\rrbracket$
provided we add the propositional resizing axiom and the Mahlo--Universe axiom needed for
Replacement.
\end{theorem}
\begin{proof}[Sketch]
Each IZF set is interpreted as a well-founded $\mathbb{V}$-structure
(Aczel's V) inside HoTT. Full Separation translates to subtype formation under
propositional resizing. Powerset translates to $A \to \mathrm{hProp}$. Collection requires
the Mahlo-style closure on the universe. See HoTT Book~\cite{HoTTBook} \S10.5 for details.
\end{proof}

\subsection{Diaconescu's theorem and IZF}

\begin{theorem}[Diaconescu 1975]
\label{thm:diaconescu}
In any topos satisfying the axiom of choice, the law of excluded middle holds. In
particular, $\mathrm{IZF}+\mathrm{AC}$ proves $\mathrm{EM}$.
\end{theorem}
\begin{proof}[Sketch]
Given a proposition $P$, consider the set $A = \{x\in\{0,1\}: x=0 \vee P\}$ and
$B = \{x\in\{0,1\}: x=1\vee P\}$. By Pairing both are sets, and their union is
$\{0,1\}$. AC gives a choice function on the family of these two non-empty sets,
producing $a\in A, b\in B$. Now either $a=b$, in which case $P$ holds (decidably), or
$a\neq b$, in which case $\neg P$ holds. Either way $P\vee\neg P$.
\end{proof}

\begin{corollary}
$\mathrm{IZF}+\mathrm{AC}$ is equivalent to $\mathrm{ZFC}$. Hence the constructive
content of IZF is genuinely tied to the absence of AC.
\end{corollary}

\subsection{Aczel's CZF and predicativity}

A subtle but important alternative to IZF is Aczel's CZF (Constructive
Zermelo--Fraenkel), formulated in~\cite{Aczel1978}. CZF differs from IZF in two crucial
ways:
\begin{itemize}
\item \emph{Restricted Separation}: the Separation schema is restricted to bounded
formulas (formulas all of whose quantifiers are restricted to a set).
\item \emph{Subset Collection (Fullness)}: replaces Powerset.
\end{itemize}
The result is a predicative theory, in the sense of Feferman: CZF can be interpreted
in Martin-L\"of type theory (the sets-in-types interpretation~\cite{Aczel1978}), and
hence has the proof-theoretic strength of MLTT plus a single Mahlo-like inaccessible.

\begin{remark}
The IZF/CZF distinction matters for the comparison with HoTT: CZF is closer to HoTT in
spirit (predicative, type-theoretic), whereas IZF is closer to ZF in spirit
(impredicative, set-theoretic). Both can be interpreted in HoTT, but CZF more
naturally.
\end{remark}

\section{FOLDS: First-Order Logic with Dependent Sorts}
\label{sec:folds}

\subsection{Dependent signatures}

\begin{definition}[FOLDS signature, after Makkai~\cite{Makkai1995}]
\label{def:foldssig}
A \emph{FOLDS signature} is a small category $\mathbf{L}$ that is:
\begin{itemize}
\item \emph{One-way}: there are no non-trivial endomorphisms or antiparallel parallel
arrows.
\item \emph{Direct}: every object has finitely many objects above it, and the order on
``levels'' (height in the underlying graph) is well-founded.
\item Equipped with a designated set of \emph{relation symbols}: distinguished objects
intended to be interpreted as truth-valued (i.e.\ subterminal).
\end{itemize}
A \emph{kind} is a non-relation object; a \emph{relation} is one of the designated objects.
\end{definition}

\begin{example}[FOLDS signature for categories]
\label{ex:foldscat}
The signature $\mathbf{L}_{\mathrm{Cat}}$ has:
\begin{itemize}
\item Kind $O$ (objects).
\item Kind $A$ (arrows), with two morphisms $s,t: A\to O$ (source, target). The
dependency $A\to O\times O$ records that an arrow has a source and target object.
\item Relation $T(g,h,k)$ on composable triples $(g, h, k:A)$ (where $t(g)=s(h)$ and the
common boundary types match) expressing ``$k$ is the composite $h\circ g$''. This
relation captures composition data without committing to a function symbol.
\item Relation $I(f)$ on $A$ expressing ``$f$ is an identity arrow''. We list this
relation explicitly because in FOLDS one cannot pick out identity arrows by equality
with a chosen constant; instead, identities are designated via this relation, which is
required to satisfy the laws $\forall f.\, I(f)\to T(f,g,g)$ and $T(g,f,g)$ for parallel
$g$.
\item Relation $E(f,g)$ on parallel arrows expressing arrow-equality. In FOLDS the
equality predicate on each kind is itself a designated relation symbol, not a primitive
of the logic; this is what guarantees isomorphism-invariance of FOLDS sentences.
\end{itemize}
A model is a category in the usual sense, where the three relations $T, I, E$ have
distinct and complementary roles: $T$ records composition data, $I$ singles out the
identity arrows, and $E$ provides arrow-equality. Together they suffice to express the
category axioms (associativity, unit laws, identity laws) as FOLDS sentences.
\end{example}

\subsection{FOLDS-equivalence}

\begin{definition}[FOLDS-equivalence, Makkai~\cite{Makkai1995}]
\label{def:foldsequiv}
Two $\mathbf{L}$-structures $M, N$ are \emph{FOLDS-equivalent} (written $M\simeq_F N$)
iff there exists a span of $\mathbf{L}$-structure homomorphisms
\[
M \xleftarrow{p} P \xrightarrow{q} N
\]
where both $p$ and $q$ are \emph{very surjective}: surjective on every kind, and on every
relation symbol $R$, the induced map $R^P \to R^M\times_{(\text{kinds})} \cdots$ is
surjective.
\end{definition}

\begin{remark}
For the signature $\mathbf{L}_{\mathrm{Cat}}$, $M\simeq_F N$ coincides with the
classical notion of equivalence of categories $M\simeq N$ (essentially surjective and
fully faithful). For higher signatures (e.g.\ bicategories), FOLDS-equivalence
correctly generalises bi-equivalence; for $n$-categories, $n$-equivalence; and so on for
$\infty$-categorical signatures.
\end{remark}

\subsection{Makkai's Invariance Theorem}

\begin{theorem}[Makkai's Invariance Theorem~\cite{Makkai1995}]
\label{thm:invariance}
For every FOLDS signature $\mathbf{L}$ and every FOLDS sentence $\varphi$ in
$\mathcal{L}(\mathbf{L})$, and any two $\mathbf{L}$-structures $M,N$:
\[
M\simeq_F N \;\Longrightarrow\; (M\models \varphi \Leftrightarrow N\models \varphi).
\]
\end{theorem}
\begin{proof}[Sketch]
By induction on $\varphi$. Atomic formulas are preserved by very surjective
homomorphisms, both ways, by the surjectivity on relations. Connectives and the unique
quantifier $\exists x: K$ (where $K$ is a kind) preserve preservation under very
surjective spans because every element on each kind in $M$ (resp.\ $N$) is hit by some
element of $P$. Full induction is given in
\cite[Theorem 4.4]{Makkai1995}; an alternative homotopy-theoretic proof appears in
Henry's work on the Isomorphism Property.
\end{proof}

\begin{corollary}
\label{cor:invariance}
No FOLDS sentence can distinguish two equivalent categories. In particular, the predicate
``$X$ is the $\in$-rank-zero element of the universe'' is \emph{not} expressible in any
FOLDS signature for ZF.
\end{corollary}

\subsection{FOLDS as a fragment of dependent type theory}

Palmgren~\cite{Palmgren2016} showed FOLDS embeds into intensional Martin-L\"of type
theory. The dependence in a FOLDS signature is captured by the dependence of types in
MLTT; relations become $(-1)$-types (mere propositions).

\begin{theorem}[Palmgren~\cite{Palmgren2016}]
\label{thm:palmgren}
There is a faithful translation $\mathcal{T}$ from FOLDS signatures over a category
$\mathbf{L}$ to dependent type theories with explicit substitutions
$\mathrm{CwF}(\mathbf{L})$, such that:
\begin{enumerate}
\item Every FOLDS structure $M$ on $\mathbf{L}$ corresponds to a model of
$\mathcal{T}(\mathbf{L})$.
\item FOLDS-equivalence corresponds to equivalence of models of $\mathcal{T}(\mathbf{L})$.
\end{enumerate}
\end{theorem}

\subsection{Worked FOLDS proofs}

\begin{example}[Categorical equivalence is FOLDS-equivalence]
\label{ex:catequivfolds}
Take $\mathbf{L} = \mathbf{L}_{\mathrm{Cat}}$ and let $F:M\to N$ be a functor. Construct
the cograph (mapping cylinder) $C(F)$: objects are
$\mathrm{Ob}(M)\sqcup\mathrm{Ob}(N)$, with hom-sets given by
\[
C(F)(x,y) := \begin{cases}
M(x,y) & x,y\in\mathrm{Ob}(M)\\
N(x,y) & x,y\in\mathrm{Ob}(N)\\
N(F(x), y) & x\in M, y\in N\\
\emptyset & \text{otherwise.}
\end{cases}
\]
Then both inclusions $M\hookrightarrow C(F)$ and $N\hookrightarrow C(F)$ are
homomorphisms of $\mathbf{L}_{\mathrm{Cat}}$-structures. The inclusion of $M$ is
\emph{very surjective} iff $F$ is essentially surjective and full. The inclusion of $N$
is very surjective iff $F$ is faithful.  Combining, $F$ is an equivalence iff there is a
zigzag of cograph cospans connecting $M$ and $N$ via very surjective maps --- i.e.\
$M\simeq_F N$ in the FOLDS sense.
\end{example}

\begin{proposition}[Skeletality is not FOLDS-expressible]
\label{prop:skeleton-not-folds}
The predicate ``$\mathcal{C}$ is a skeletal category'' (every object is the unique
representative of its isomorphism class) is not expressible by any FOLDS sentence.
\end{proposition}
\begin{proof}
Suppose $\varphi$ is a FOLDS sentence in $\mathbf{L}_{\mathrm{Cat}}$ that holds in
$\mathcal{C}$ iff $\mathcal{C}$ is skeletal. Take a non-skeletal $\mathcal{C}_0$ and a
skeleton $\mathcal{C}_1\hookrightarrow \mathcal{C}_0$. The inclusion is an
equivalence (essentially surjective and fully faithful), hence the cograph construction
yields $\mathcal{C}_0 \simeq_F \mathcal{C}_1$. By Makkai's invariance theorem
(\cref{thm:invariance}), $\mathcal{C}_0\models \varphi \Leftrightarrow \mathcal{C}_1
\models\varphi$. But $\mathcal{C}_1$ is skeletal and $\mathcal{C}_0$ is not, so $\varphi$
cannot distinguish them. Contradiction.
\end{proof}

\begin{remark}
\Cref{prop:skeleton-not-folds} is paradigmatic: any predicate referring to specific
\emph{representatives} of equivalence classes (rather than equivalence-classes themselves)
fails to be FOLDS-expressible. This is the syntactic mechanism by which FOLDS enforces
the equivalence principle: ``junk'' predicates cannot even be \emph{stated}.
\end{remark}

\subsection{Higher signatures: bicategories and \texorpdfstring{$\infty$}{infinity}-categories}

The signature $\mathbf{L}_{\mathrm{Cat}}$ for ordinary categories generalises naturally
to higher signatures:
\begin{itemize}
\item $\mathbf{L}_{\mathrm{Bicat}}$: kinds $O$, $A$, $C$ (cells), with
$s,t:A\to O$, $s,t:C\to A$, and relations for vertical/horizontal composition,
identities, and 2-cell equality. FOLDS-equivalence on $\mathbf{L}_{\mathrm{Bicat}}$
is bi-equivalence of bicategories.
\item $\mathbf{L}_{(\infty,1)}$: an infinite tower of kinds $A_n$ for $n$-cells, with
boundary maps and a single ``hcomp'' relation per dimension. FOLDS-equivalence
corresponds to Joyal-style $\infty$-equivalence (DK-equivalence).
\end{itemize}
Szumi\l{}o~\cite{Szumilo2019} extends this to a hierarchy of $n$-FOLDS logics whose
syntax converges to the syntax of HoTT.

\section{HoTT: The unifying frame}
\label{sec:hott}

\subsection{Univalence as the equivalence-invariance axiom}

Univalence (Voevodsky 2010, HoTT Book~\cite{HoTTBook} \S2.10) is the axiom:
\[
\mathsf{ua}: (A\simeq B) \;\xrightarrow{\sim}\; (A=_{\mathcal{U}} B)
\]
saying that the canonical map ``$\mathrm{idtoeqv}$'' from identifications between types
to equivalences between them is itself an equivalence. In short: \emph{equivalent types
are identical}. The Structure Identity Principle (SIP) generalises this from raw types to
typed structures: equivalent structures are identical.

\subsection{Which structural principles become theorems}

The following table summarises what passes from axiom-status (in ETCS, IZF, or as a
metatheorem in FOLDS) to theorem-status in HoTT (Section 1 column reproduced for
convenience):
\begin{itemize}
\item \textbf{NNO} is a theorem in HoTT --- the inductive type $\mathbb{N}$ is built into
the type theory and contractibility of the type of NNO structures is provable
(Paper~V Theorem~4.4 of our series).
\item \textbf{Function extensionality} is a theorem in HoTT given univalence: Voevodsky's
proof shows $\mathrm{UA}\Rightarrow \mathrm{FE}$.
\item \textbf{Propositional extensionality} is a theorem in HoTT given univalence,
restricted to mere propositions.
\item \textbf{Structure Identity Principle} is a theorem in HoTT given univalence
(see Paper~VI Theorems 10.3--10.4 of our series).
\item \textbf{FOLDS-equivalence-implies-identity} is a theorem in HoTT given univalence
(this is the Univalence Principle of Ahrens--North--Shulman--Tsementzis 2019).
\end{itemize}

We now make this last statement precise.

\subsection{The Univalence Principle (ANST 2019)}

\begin{theorem}[Univalence Principle, ANST~\cite{ANST2019}]
\label{thm:up}
For any FOLDS signature $\mathbf{L}$ and any two $\mathbf{L}$-structures $M, N$ in HoTT,
\[
(M\simeq_F N)\;\simeq\;(M=_{\mathrm{Str}(\mathbf{L})} N).
\]
That is, FOLDS-equivalence is identification in the type of $\mathbf{L}$-structures.
\end{theorem}
\begin{proof}[Sketch]
The dependent-sort structure of $\mathbf{L}$ is encoded as a Reedy diagram of types in
HoTT. Each relation symbol becomes a $(-1)$-truncated type, each kind a type of the
appropriate truncation level. By induction on the Reedy structure, FOLDS-equivalence
unfolds into the data of an equivalence at each level, which by repeated application of
the Structure Identity Principle (a consequence of univalence) gives an identification.
The full proof requires careful management of the level structure of $\mathbf{L}$ and the
Reedy fibrant replacement; see ANST.
\end{proof}

\begin{remark}
\Cref{thm:up} subsumes Makkai's Invariance Theorem: indeed if $M\simeq_F N$ then
$M=N$ in HoTT, hence any predicate is preserved (by transport along the identification).
But \cref{thm:up} is stronger: it gives an \emph{equivalence}, not just an implication.
The strength comes from univalence.
\end{remark}

\subsection{Internalising ETCS in HoTT}

We now sketch how ETCS is recovered as a theorem inside HoTT, given univalence and
sufficient choice/EM where needed.

\begin{theorem}[Sets in HoTT form a model of ETCS-AC]
\label{thm:setsmodel}
Within HoTT, define $\mathrm{Set}$ to be the type of $0$-truncated types (sets in the
sense of HoTT \S3.1):
\[
\mathrm{Set} := \Sigma_{X:\mathcal{U}}\,\mathrm{isSet}(X).
\]
Then $\mathrm{Set}$, with the obvious morphism structure (functions between sets) and
identity and composition, is an elementary topos with NNO. Adding propositional resizing and
the (provably independent) axioms EM and AC produces a model of ETCS.
\end{theorem}
\begin{proof}[Sketch]
\textbf{Finite limits}: Sigma-types and product types are sets when their components are
sets; equalisers are $\Sigma_{x:A} (f(x) = g(x))$, also a set.
\textbf{Cartesian closure}: function types $A\to B$ are sets when $B$ is a set (by FE).
\textbf{Subobject classifier}: $\Omega := \mathrm{hProp}$, the type of mere propositions;
truth is the standard map.
\textbf{NNO}: the inductive type $\mathbb{N}$.
\textbf{Well-pointedness}: function extensionality.
\textbf{AC}: axiom (independent in HoTT).
\textbf{Two-valuedness}: requires EM (axiom).
\textbf{Non-degeneracy}: $\mathbf{0}\not\simeq\mathbf{1}$ holds in any non-trivial model.
\end{proof}

\subsection{Internalising IZF in HoTT}

Aczel's $\mathbb{V}$-construction inside HoTT:

\begin{definition}[The HoTT cumulative hierarchy]
\label{def:vhott}
Define the higher inductive type $\mathbb{V}$ with:
\begin{itemize}
\item Constructor $\mathrm{set}: (A:\mathcal{U}) \to (A\to \mathbb{V})\to \mathbb{V}$.
\item Path constructor: extensionality, $\mathrm{set}(A,f) =_{\mathbb{V}} \mathrm{set}(B,g)$
iff there is a bisimulation between the trees $(A,f)$ and $(B,g)$.
\item $0$-truncation.
\end{itemize}
\end{definition}

\begin{theorem}[HoTT Book \S10.5]
\label{thm:vmodelizf}
$\mathbb{V}$ is a model of IZF in the following sense: every IZF axiom (Extensionality,
Pairing, Union, Powerset, Infinity, Collection, Full Separation, $\in$-Induction)
translates to a provable HoTT statement about $\mathbb{V}$, given propositional resizing.
\end{theorem}

\subsection{Internalising FOLDS in HoTT via the Univalence Principle}

The key observation: FOLDS signatures, qua small categories, are themselves objects of
HoTT. FOLDS-structures qua functors are objects of HoTT. The Univalence Principle
(\cref{thm:up}) is the statement that the type of $\mathbf{L}$-structures has the
expected identity type. So FOLDS is not just \emph{representable} in HoTT; the FOLDS
metatheorem (Makkai's Invariance Theorem) becomes an internal HoTT theorem.

\section{Models}
\label{sec:models}

\subsection{ETCS in toposes}

A model of ETCS is a well-pointed elementary topos with NNO and AC. The category
$\mathrm{Set}$ in ZFC is the prototype. Other models include:
\begin{itemize}
\item $\mathrm{Set}_{\mathrm{ZF+AC}}$: classical sets in ZF + AC.
\item $\mathrm{Set}_{\mathrm{IZF+AC}}$: intuitionistic sets in IZF + AC (where AC implies
EM via Diaconescu, hence equivalent to ZFC).
\item $\mathrm{FinSet}$: finite sets satisfy A1--A6, fail A4 (no NNO), so not ETCS.
\item Forcing extensions of $\mathrm{Set}$: ETCS+CH and ETCS+$\neg$CH.
\end{itemize}

\subsection{IZF and topos-valued models}

Friedman--Aczel construction: in any cocomplete elementary topos $\mathcal{E}$ with NNO
and a Mahlo-style universe object, the internal $\mathbb{V}$-structure (the
``cumulative hierarchy'' inside $\mathcal{E}$) gives a model of IZF.

\begin{theorem}[Friedman--Aczel--Joyal]
The internal hierarchy $V_{\mathcal{E}}$ in a cocomplete topos with universe object
satisfies all axioms of IZF.
\end{theorem}

The proof reduces each axiom (Extensionality, Pairing, Union, Powerset, Infinity,
Collection, Full Separation, $\in$-Induction) to a property of $\mathcal{E}$ that holds
by definition of cocomplete elementary topos with universe.

\subsection{FOLDS in fibrations}

A FOLDS signature $\mathbf{L}$ corresponds to a discrete Conduch\'e fibration over
$\mathbf{L}^{\mathrm{op}}$. Models of $\mathbf{L}$ are categories of elements/sections of
this fibration.

\begin{theorem}[Makkai~\cite{Makkai1995}]
\label{thm:foldsfib}
The category $\mathrm{Mod}(\mathbf{L})$ of $\mathbf{L}$-structures and homomorphisms is
the slice $\mathrm{Set}/\mathbf{L}^{\mathrm{op}}$ when $\mathbf{L}$ is direct (with
appropriate naturality), with FOLDS-equivalence being the localisation along
very-surjective maps.
\end{theorem}

\subsection{HoTT in \texorpdfstring{$(\infty,1)$}{(infinity,1)}-toposes}

\begin{theorem}[Shulman~\cite{Shulman2019}]
\label{thm:shulman}
Every Grothendieck $(\infty,1)$-topos models HoTT with univalent universes.
\end{theorem}
\begin{proof}[Sketch]
Shulman's strategy: use simplicial categories of fibrant types to build an interpretation
of MLTT in any presentable $(\infty,1)$-topos $\mathcal{X}$, then use
Reedy fibrant replacement and the local universe construction (Lumsdaine--Warren) to
produce a univalent universe object.
\end{proof}

\subsection{The square of interpretations}

\begin{center}
\begin{tikzcd}[column sep=large, row sep=large]
\mathrm{ETCS} \arrow[r, "\text{McLarty}"] \arrow[d, "\text{Joyal}"'] &
\mathrm{BZ}+\mathrm{R} \arrow[d, "\text{Friedman}"] \\
\mathrm{Topos+NNO+AC} \arrow[r, "\text{int.\ lang.}"'] &
\mathrm{IZF}+\mathrm{AC} \arrow[r, leftrightarrow, "\text{Diaconescu}"]&\mathrm{ZFC}
\end{tikzcd}
\end{center}

with FOLDS sitting orthogonal to all three (FOLDS is a logic, not a set theory; it
specifies which sentences are admissible). HoTT sits ``above'' the diagram: any model of
ETCS, IZF, or FOLDS yields a model of HoTT after passage through the
$(\infty,1)$-topos completion.

\subsection{Reverse-mathematical strengths}

We summarise the proof-theoretic strength of each system:

\begin{center}
\begin{tabular}{lll}
\toprule
\textbf{System} & \textbf{Strength (consistency)} & \textbf{Reference} \\
\midrule
ETCS & ZFC $-$ Replacement & McLarty 2004 \\
ETCS+R (ETCS with Replacement) & ZFC & McLarty 2004 \\
IZF & ZF & Friedman 1973 \\
IZF+AC & ZFC & Diaconescu 1975 \\
CZF & MLTT $+$ universes & Aczel 1978 \\
CZF $+$ Mahlo & ZFC & Rathjen \\
HoTT (with UA) & ZFC $+$ inaccessibles & Shulman 2019 \\
HoTT (without UA) & MLTT & MLTT proof theory \\
FOLDS (as logic) & --- & --- \\
\bottomrule
\end{tabular}
\end{center}

The interesting feature: ETCS minus Replacement is strictly weaker than ZFC; recovering
full ZFC strength requires a Replacement schema, in either material or structural form.
HoTT with univalence and propositional resizing matches ZFC + countably many
inaccessibles (Shulman~\cite{Shulman2019}), since each universe behaves as a Grothendieck
universe.

\section{The univalence boundary}
\label{sec:boundary}

\subsection{Definition}

\begin{definition}[Univalence boundary]
\label{def:boundary}
A structural principle $P$ \emph{lies on the far side of the univalence boundary} if
$P$ holds in HoTT~+~UA but not in HoTT alone (i.e.\ in MLTT plus the assumption that
identity types are well-behaved but without the univalence axiom). A principle lies on
the \emph{near side} if it holds in HoTT without UA. We say $P$ \emph{straddles} the
boundary if it is independent of UA in HoTT.
\end{definition}

\subsection{Far-side principles}

\begin{theorem}[Far-side: SIP needs UA]
\label{thm:sipneedsua}
The Structure Identity Principle (SIP) is provably equivalent to univalence in HoTT.
\end{theorem}
\begin{proof}[Sketch]
$\mathrm{UA}\Rightarrow \mathrm{SIP}$: the standard structure-identity proof (HoTT
Book \S9.8) constructs the identification of structures using univalence applied to the
underlying types and transport for the operations.
$\mathrm{SIP}\Rightarrow \mathrm{UA}$: instantiate SIP at the trivial signature
(structure-free) to extract univalence as the special case where the only operation is
the identity.
\end{proof}

\begin{theorem}[Far-side: FOLDS-equivalence-implies-identity needs UA]
\label{thm:foldsneedsua}
The principle ``$M\simeq_F N \Rightarrow M=_{\mathrm{Str}(\mathbf{L})}N$'' for arbitrary
FOLDS signatures $\mathbf{L}$ is equivalent to univalence (over MLTT with FE and PE).
\end{theorem}
\begin{proof}[Sketch]
$\mathrm{UA}\Rightarrow$ statement: this is \cref{thm:up}.
Statement $\Rightarrow \mathrm{UA}$: take $\mathbf{L}$ to be the signature with one kind
and no relations. A model is just a type. FOLDS-equivalence on this signature is type
equivalence. The conclusion that equivalence is identification is exactly UA.
\end{proof}

\subsection{Near-side principles}

\begin{theorem}[Near-side: NNO existence is near]
\label{thm:nnonear}
The existence of an NNO in HoTT (as the inductive type $\mathbb{N}$) does not require UA.
\end{theorem}
\begin{proof}
$\mathbb{N}$ is generated by the constructors $0:\mathbb{N}$ and $s:\mathbb{N}\to
\mathbb{N}$ together with the universal property given by the recursor; this is the
elimination rule of MLTT inductive types and does not invoke UA.
\end{proof}

\begin{theorem}[Near-side: function extensionality is near, given univalence is far]
\label{thm:fefar}
Function extensionality (FE) holds in HoTT~+~UA but is independent of MLTT alone.
However, FE itself does not imply UA.
\end{theorem}
\begin{proof}[Sketch]
$\mathrm{UA}\Rightarrow \mathrm{FE}$: Voevodsky's proof.
$\mathrm{FE}\not\Rightarrow \mathrm{UA}$: the simplicial set model satisfies FE but
without univalent universes (univalence requires more structure).
\end{proof}

\subsection{Straddlers}

\begin{theorem}[Straddler: EM, AC are independent]
\label{thm:straddlers}
The law of excluded middle EM and the axiom of choice AC are independent of UA in HoTT.
\end{theorem}
\begin{proof}[Sketch]
EM: the classical model (sets) satisfies EM and UA; the topological model satisfies UA
but not EM. Similarly AC.
\end{proof}

\subsection{The boundary as a square}

\begin{center}
\begin{tabular}{lcc}
\toprule
& \textbf{Without UA} & \textbf{With UA} \\
\midrule
\textbf{NNO} & yes & yes \\
\textbf{FE} & independent & yes \\
\textbf{PE} & independent & yes \\
\textbf{SIP} & no & yes \\
\textbf{FOLDS-eq=id} & no & yes \\
\textbf{EM} & independent & independent \\
\textbf{AC} & independent & independent \\
\textbf{Replacement} & yes (small) & yes (small) \\
\bottomrule
\end{tabular}
\end{center}

\subsection{Diagnostic principles}

Given an arbitrary structural principle $P$, how can one test whether $P$ lies on the
near or far side of the boundary? We propose the following heuristic.

\begin{proposition}[Diagnostic for far-side principles]
\label{prop:diagnostic}
A principle $P$ lies on the far side of the univalence boundary if and only if $P$
asserts that two equivalent structures are \emph{equal}, where equality is in the
sense of HoTT's identity type.
\end{proposition}
\begin{proof}[Heuristic argument]
The forward direction is clear: such principles are precisely instances of the SIP
schema (or the more general Univalence Principle), and we showed in
\cref{thm:sipneedsua,thm:foldsneedsua} that these need UA.
The reverse direction relies on a meta-theoretic observation: principles that are not
about identification are about \emph{existence} or \emph{construction}, and these are
unaffected by UA's content (which is about identity). A detailed proof requires
formalisation of the schema, deferred to~\cite{ANST2019}.
\end{proof}

\subsection{What is on the boundary?}

A principle is \emph{on the boundary} if it is equivalent to UA in HoTT. Examples:
\begin{itemize}
\item SIP (\cref{thm:sipneedsua}).
\item Univalence Principle for FOLDS (\cref{thm:foldsneedsua}).
\item Function extensionality + propositional extensionality + Voevodsky's
``equivalence-induction'' rule (a known characterisation of UA, due to Voevodsky).
\item Equivalence of $\Pi$-types respecting equivalence of arguments (a partial
characterisation).
\end{itemize}

Principles strictly weaker than UA (i.e.\ near side or straddle):
\begin{itemize}
\item FE alone is strictly weaker than UA (the simplicial model with non-univalent
universes satisfies FE).
\item PE alone is strictly weaker than UA.
\item NNO existence (theorem of MLTT, no UA needed).
\end{itemize}

\section{Open problems}
\label{sec:open}

\subsection{Problem 1: Contractibility of the type of ETCS structures}

\begin{conjecture}
\label{conj:etcscontract}
In HoTT~+~UA, the type
\[
\mathrm{ETCS}\text{-Str} := \Sigma_{\mathcal{S}:\mathrm{Cat}_\infty}\,\Pi_{i=1}^{8}
\mathrm{Ax}_i(\mathcal{S})
\]
is contractible up to a fixed choice of universe parameter.
\end{conjecture}

\noindent A proof would mirror Paper~V Theorem~4.4 (contractibility of the NNO type) but
at the level of the entire categorical universe. Tools needed: SIP for $\infty$-toposes,
internal $\infty$-Yoneda (Riehl--Shulman directed type theory).

\subsection{Problem 2: Cubical Agda formalisation of FOLDS}

\begin{conjecture}
\label{conj:cubicalfolds}
There is a Cubical Agda library implementing FOLDS signatures, FOLDS-equivalence, and
$n$-level FOLDS (Szumi\l{}o 2019), and proving \cref{thm:up} formally for $n\le 2$
(equivalence of FOLDS-equivalence and HoTT-identification at $n=0,1,2$).
\end{conjecture}

\noindent Partial work exists in the agda-categories library and in Tsementzis's
formalisation of the SIP. Full FOLDS in Cubical Agda would close the gap between Makkai
1995 and ANST 2019 syntactically.

\subsection{Problem 3: Identifying IZF axioms with HoTT principles}

\begin{conjecture}
\label{conj:izfhott}
There is a precise correspondence between the IZF axioms and HoTT principles of the
following form:
\begin{itemize}
\item Extensionality $\leftrightarrow$ FE + PE.
\item Pairing/Union $\leftrightarrow$ Sigma/Coproduct.
\item Powerset $\leftrightarrow$ $A\to\mathrm{hProp}$ (with propositional resizing).
\item Collection $\leftrightarrow$ Mahlo universe axiom.
\item $\in$-Induction $\leftrightarrow$ accessibility predicate.
\item Full Separation $\leftrightarrow$ subtype formation under propositional resizing.
\end{itemize}
\end{conjecture}

\noindent Partial work: HoTT Book \S10.5; Awodey--Warren; Rathjen--Tupailo on CZF.

\section{Discussion}
\label{sec:discussion}

\subsection{What does ``foundation'' mean today?}

The three-way comparison reveals that ``foundations'' has split into three orthogonal
axes:
\begin{itemize}
\item \textbf{Material vs structural}: ZFC (material) vs ETCS (structural).
\item \textbf{Classical vs constructive}: ZFC vs IZF/CZF.
\item \textbf{0-categorical vs higher}: ETCS (1-categorical) vs HoTT ($\infty$-categorical).
\end{itemize}
FOLDS is orthogonal to all three: it is a \emph{logic}, applicable in any of the
combinations.

\subsection{Practical recommendations}

For the working mathematician:
\begin{itemize}
\item \textbf{Pure category theory, classical}: ETCS+AC.
\item \textbf{Pure category theory, constructive}: well-pointed Heyting topos.
\item \textbf{Higher category theory, classical}: HoTT+UA+AC (= simplicial sets model).
\item \textbf{Higher category theory, constructive}: HoTT+UA (= cubical model).
\item \textbf{Pure logic, isomorphism-invariant}: FOLDS over HoTT (Univalence Principle).
\item \textbf{Reverse-mathematical strength}: stay with ZFC variants.
\end{itemize}

\subsection{Limitations}

This paper does not prove the full Univalence Principle; we report the result of ANST
2019. We also do not address the predicativity issue: HoTT proper is predicative, IZF is
impredicative, ETCS sits in between. A precise predicativity-comparative analysis is
left to future work. Finally, the ``square of interpretations'' is not commutative on
the nose; the diagram commutes up to bi-interpretation, a notion we have not fully
formalised here.

\section{Worked case study: monoid theory across the four foundations}
\label{sec:casestudy}

To make the comparison concrete, we trace a single mathematical concept --- the theory
of monoids --- through ETCS, IZF, FOLDS, and HoTT. This illustrates how
``structural-ness'' progresses from convention (ETCS) to enforced syntax (FOLDS) to
internal theorem (HoTT+UA).

\subsection{Monoids in ZFC and IZF}

In ZFC/IZF, a monoid is a 4-tuple $(M, \cdot, e, \text{laws})$ where $M$ is a set, $\cdot:
M\times M\to M$, $e\in M$, and the associativity and unit laws hold. The encoding of the
4-tuple uses Kuratowski pairs. Two monoids are \emph{equal} iff they have the same
underlying $\in$-tree --- a notion that depends on the choice of pair encoding. Monoid
isomorphism is a derived notion: a bijection $\phi: M\to N$ commuting with operations.
Junk theorems (e.g.\ ``the underlying set of $\mathbb{Z}$ contains the empty set as an
element'') depend on encoding.

\subsection{Monoids in ETCS}

In ETCS, a monoid is an object $M$ together with morphisms $\cdot: M\times M\to M$,
$e: 1\to M$, and the equational axioms expressed using composition and the universal
property of products. There is no encoding choice for $\cdot$ or $e$. Two monoids are
isomorphic iff there is an invertible morphism between them respecting the operations.
The category $\mathrm{Mon}(\mathrm{ETCS})$ of monoids in ETCS is well-defined, and is the
category of models of the algebraic theory of monoids.

\subsection{Monoids in FOLDS}

A FOLDS signature for monoids:
\begin{itemize}
\item Kind $M$ (elements).
\item Relation $E(x,y)$ (equality, parallel arrows from triangulation).
\item Relation $P(x,y,z)$ (multiplication: $z = x\cdot y$).
\item Relation $U(x)$ (unit: $x = e$).
\end{itemize}
Axioms:
\begin{itemize}
\item Associativity:
\[
\begin{aligned}
\forall a,b,c,xy,yz,abc,abc': M.\;
& P(a,b,xy) \wedge P(xy,c,abc) \wedge P(b,c,yz) \\
& {} \wedge P(a,yz,abc') \;\Rightarrow\; E(abc, abc').
\end{aligned}
\]
\item Left unit: $\forall e, x, ex: M.\, U(e) \wedge P(e,x,ex) \Rightarrow E(ex, x)$.
\item Right unit: dual.
\end{itemize}
A FOLDS-equivalence between two monoids $M, N$ is a span via very surjective
homomorphisms, which here coincides with the usual monoid isomorphism. Makkai's
invariance theorem guarantees: any FOLDS sentence holding in $M$ also holds in $N$.
There is no way to express ``$M$ is the monoid whose elements are von Neumann ordinals
$\le 5$'' in FOLDS, since this is non-isomorphism-invariant.

\subsection{Monoids in HoTT}

In HoTT, a monoid is a Sigma-type:
\[
\mathrm{Monoid} := \Sigma_{M:\mathrm{Set}}\Sigma_{\cdot: M\to M\to M}\Sigma_{e:M}
\bigl(\mathrm{Assoc}(\cdot)\times \mathrm{LUnit}(\cdot,e)\times \mathrm{RUnit}(\cdot,e)\bigr).
\]
The components are: a set, a binary operation, an identity element, and proofs (mere
propositions) of the laws.

\begin{theorem}[Structure Identity Principle for monoids]
\label{thm:sipmonoid}
For two monoids $M, N: \mathrm{Monoid}$:
\[
(M =_{\mathrm{Monoid}} N) \;\simeq\; \mathrm{MonoidIso}(M, N)
\]
where $\mathrm{MonoidIso}(M,N)$ is the type of monoid isomorphisms.
\end{theorem}
\begin{proof}[Sketch]
By Sigma-induction on $M, N$ and SIP for each component: the sets are identified by UA,
the operations and identity are transported, the laws are mere propositions and so
identified automatically. The full proof is a routine application of HoTT Book \S9.8.
\end{proof}

\subsection{Comparison}

\begin{center}
\begin{tabular}{llll}
\toprule
& \textbf{Equality} & \textbf{Iso} & \textbf{SIP-statement} \\
\midrule
ZFC & encoding-dependent & derived & metatheorem (sociological) \\
IZF & encoding-dependent & derived & metatheorem \\
ETCS & up to unique iso & primitive & metatheorem (categorical) \\
FOLDS & up to FOLDS-eq & primitive & syntactic (invariance theorem) \\
HoTT+UA & SIP-statement & primitive & internal theorem \\
\bottomrule
\end{tabular}
\end{center}

The rightmost column is the key: only HoTT+UA makes the SIP an internal theorem, expressed
in the same language as the structures themselves. ETCS makes it a category-theoretic
metatheorem (true of every property of the category of sets). FOLDS makes it a
syntactic invariance theorem (true of every well-formed FOLDS sentence). ZFC and IZF
make it a sociological convention. Each escalation reduces the ``trust assumption'' needed.

\section{Practical foundational choice}
\label{sec:practical}

\subsection{Which foundation for which task?}

The foundational pluralism we have surveyed is not just a curiosity; it has practical
consequences for formalisation projects.

\begin{itemize}
\item \textbf{Number theory and analysis}: ZFC variants are best, due to the deep classical
results requiring AC and EM. Lean/Mathlib uses ZFC (via the type-theoretic
formalisation of $\mathsf{ZFSet}$ inside Lean).
\item \textbf{Category theory and topos theory}: ETCS or HoTT+UA. ETCS is sufficient for
1-categorical structuralism; HoTT+UA is needed for higher-categorical structuralism (the
theory of $(\infty,1)$-categories).
\item \textbf{Synthetic differential geometry, synthetic homotopy theory}: HoTT, naturally,
since the synthetic methods rely on the type-theoretic semantics.
\item \textbf{Constructive analysis and computer-verified mathematics}: CZF or HoTT
without UA, where computational interpretation is crucial.
\item \textbf{Foundations of foundations} (model theory of foundations themselves): IZF
and topos theory provide the best general-purpose framework.
\item \textbf{Pedagogy}: ETCS is arguably the cleanest entry point, since it explains
``what set theory is really doing'' without the encoding noise.
\end{itemize}

\subsection{Foundational consensus or pluralism?}

A plausible reading of the present landscape is that no single foundation dominates;
each excels in a particular domain. This contrasts with the 20th century, where ZFC was
dominant. We propose to call the current situation \emph{principled pluralism}: each
foundation is correct for its target domain, and translation theorems (McLarty,
Friedman--Aczel, ANST) tie them together.

\subsection{Open foundational debates}

Several debates remain unresolved as of 2026:
\begin{itemize}
\item \emph{Predicativity vs impredicativity}: HoTT proper is predicative; IZF is
impredicative. The right level of impredicativity for a unified foundation is contested.
\item \emph{Constructivity vs classicality}: HoTT defaults to constructive but is
compatible with EM; IZF is constructive; ZFC is classical; FOLDS is logic-agnostic.
\item \emph{Higher coherence}: how high should the dimensional ladder go? HoTT goes to
$\infty$; FOLDS-Szumi\l{}o goes to arbitrary $n$; ETCS stops at $1$.
\item \emph{The role of universes}: HoTT has a Russellian hierarchy; ETCS has none
(the category of all sets is the universe); IZF has none (the proper class of all sets);
FOLDS has small categories of structures.
\end{itemize}

\section{Conclusion}
\label{sec:conclusion}

We have presented ETCS, IZF, and FOLDS as three orthogonal responses to the same
underlying problem: how to build a foundation in which only structurally relevant
assertions are formulable or provable. ETCS attacks the problem at the level of
\emph{ontology} (replace sets-with-membership by sets-with-functions). IZF attacks the
problem at the level of \emph{logic} (move to intuitionistic logic, where structural
properties are easier to certify). FOLDS attacks the problem at the level of
\emph{syntax} (build a logic in which only invariant predicates are expressible).

HoTT---in particular HoTT~+~UA---synthesises all three. The Structure Identity Principle
recapitulates the structural ontology of ETCS at the level of types. Constructive HoTT
(without EM, without AC) shares the logical environment of IZF. The Univalence Principle
recapitulates the syntactic invariance of FOLDS as a metatheorem.

The univalence boundary precisely separates: the principles that hold in MLTT alone
(the near side: NNO, induction, basic structural reasoning), and the principles requiring
univalence (the far side: SIP, FOLDS-equivalence-as-identity, the full equivalence
principle). Mapping out this boundary is, we suggest, one of the more illuminating
exercises in 2026 foundations.

\section*{Acknowledgements}

The author thanks colleagues at the YonedaAI Research Collective for discussions on
structuralism and on FOLDS-style invariance proofs. This paper continues the
foundations programme of our prior series (Papers I--VI on HoTT and structural
foundations), to which the reader is referred for the technical prerequisites.

\begin{thebibliography}{99}

\bibitem{Lawvere1964}
F.\ W.\ Lawvere,
\emph{An elementary theory of the category of sets},
Proceedings of the National Academy of Sciences 52 (1964), 1506--1511.
Reprinted with author's commentary, TAC Reprints 11 (2005).

\bibitem{LawvereRosebrugh2003}
F.\ W.\ Lawvere and R.\ Rosebrugh,
\emph{Sets for Mathematics},
Cambridge University Press, 2003.

\bibitem{McLarty2004}
C.\ McLarty,
\emph{Exploring categorical structuralism},
Philosophia Mathematica 12 (2004), 37--53.

\bibitem{Makkai1995}
M.\ Makkai,
\emph{First-order logic with dependent sorts, with applications to category theory},
Preprint, McGill University, 1995.
Available at \url{https://www.math.mcgill.ca/makkai/folds/foldsinpdf/FOLDS.pdf}.

\bibitem{Friedman1973}
H.\ Friedman,
\emph{The consistency of classical set theory relative to a set theory with intuitionistic
logic},
Journal of Symbolic Logic 38 (1973), 315--319.

\bibitem{StanfordIZFCZF}
L.\ Crosilla,
\emph{Set theory: constructive and intuitionistic ZF},
Stanford Encyclopedia of Philosophy, Spring 2024 edition, ed.\ E.\ N.\ Zalta.

\bibitem{Aczel1978}
P.\ Aczel,
\emph{The type theoretic interpretation of constructive set theory},
in Logic Colloquium '77, North-Holland (1978), 55--66.

\bibitem{Awodey2014}
S.\ Awodey,
\emph{Structuralism, invariance, and univalence},
Philosophia Mathematica 22 (2014), 1--11.

\bibitem{HoTTBook}
The Univalent Foundations Program,
\emph{Homotopy Type Theory: Univalent Foundations of Mathematics},
Institute for Advanced Study, 2013.

\bibitem{Shulman2019}
M.\ Shulman,
\emph{All $(\infty,1)$-toposes have strict univalent universes},
arXiv:1904.07004 (2019).

\bibitem{ANST2019}
B.\ Ahrens, P.\ R.\ North, M.\ Shulman, D.\ Tsementzis,
\emph{The Univalence Principle},
arXiv:2102.06275 (2019/2021).

\bibitem{Palmgren2016}
E.\ Palmgren,
\emph{Categories with families and first-order logic with dependent sorts},
arXiv:1605.01586 (2016).

\bibitem{Rasekh2018}
N.\ Rasekh,
\emph{Every elementary higher topos has a natural number object},
arXiv:1809.01734 (2018), TAC 37(13) 2021.

\bibitem{Voevodsky2010}
V.\ Voevodsky,
\emph{Univalent foundations of mathematics},
preprint, IAS Princeton (2010).

\bibitem{Hyland1982}
J.\ M.\ E.\ Hyland,
\emph{The effective topos},
in The L.\ E.\ J.\ Brouwer Centenary Symposium, North-Holland (1982), 165--216.

\bibitem{MyhillCST}
J.\ Myhill,
\emph{Constructive set theory},
Journal of Symbolic Logic 40 (1975), 347--382.

\bibitem{Joyal1984}
A.\ Joyal and I.\ Moerdijk,
\emph{Algebraic Set Theory},
LMS Lecture Note Series 220, Cambridge University Press (1995).

\bibitem{Tsementzis2017}
D.\ Tsementzis,
\emph{First-order logic with isomorphism},
arXiv:1603.03092 (2016/2017).

\bibitem{Awodey2009StructuralismInvariance}
S.\ Awodey,
\emph{From sets to types to categories to sets},
in Foundational Theories of Classical and Constructive Mathematics, Springer (2009).

\bibitem{Riehl2017}
E.\ Riehl and M.\ Shulman,
\emph{A type theory for synthetic $\infty$-categories},
Higher Structures 1 (2017), 147--224, arXiv:1705.07442.

\bibitem{Szumilo2019}
K.\ Szumi\l{}o,
\emph{Frames in cofibration categories},
Journal of Homotopy and Related Structures 14 (2019), 345--378.
See also: \emph{$n$-FOLDS and Reedy fibrant diagrams},
preprint, 2019.

\end{thebibliography}

\end{document}
